%% kapa: exam=1 round=23 year=2012 period=1 form=A subject=economics subject_ko=경제원론 seq=33 number=73
%% source: Q-Net 감정평가사 기출문제 (cst00309) artlSeq=5052578 file=raw/1008/5052578/2012년 제23회 감정평가사 제1차시험 기출문제(1교시 A형).pdf
%% conversion: kordoc 4.24.51 -> tools/gen1.py (auto); review=proofread-sonnet; flags=none
\begin{question}{73}
솔로우(R. Solow) 성장모형에서 일인당 생산함수는 $y=k^{1/2}$, 저축률은 12\%, 인구증가율은 1\%, 자본의 감가상각률은 2\%이다. 다음 설명 중 옳은 것을 모두 고른 것은? (단, $y$는 일인당 생산량, $k$는 일인당 자본량이다.)
\begin{framedtext}
ㄱ. 균제상태(steady state)에서 일인당 산출량은 4이다. \par
ㄴ. 자본소득분배율과 노동소득분배율은 같다. \par
ㄷ. 균제상태에서 황금률(golden rule)이 달성되고 있다. \par
\end{framedtext}
\begin{choices}
  \item ㄱ
  \item ㄴ
  \item ㄱ, ㄴ
  \item ㄴ, ㄷ
  \item ㄱ, ㄴ, ㄷ
\end{choices}
\end{question}
\answer{3}
